\documentclass[12pt]{article}
\usepackage[margin=1in]{geometry}
\usepackage{amsmath, amssymb, amsthm}
\usepackage{hyperref}
\usepackage{booktabs}
\usepackage{float}

\hypersetup{colorlinks=true, linkcolor=black, urlcolor=blue, citecolor=black}

\newtheorem{definition}{Definition}
\newtheorem{proposition}{Proposition}

\title{\textbf{A Phase-Alignment Coupling Model of Conscious Access:\\
Predictions from a Receiver Framework}}
\author{
    Winston Zai Lin \\
    \small Blair Academy \\
    \small \texttt{linzaiwinston413@gmail.com}
}
\date{April 2026}

\begin{document}
\maketitle

\begin{abstract}
We propose that consciousness is not generated by the body itself, but transmitted through the organised state instantiated by it. To make this hypothesis empirically testable, we formalise the body as a receiver whose state $R_i(t)$---comprising oscillation phases, synaptic weights, and neuromodulator concentrations---couples to a fixed reference phase pattern $\hat{\varphi}$ to produce individual conscious access. The coupling is given by $\varphi_i(t)=A(\tau)\cdot H(\hat{\varphi},R_i(t))$, where $A(\tau)$ decays exponentially with age. The coupling operator decomposes into layered filters; psychedelic states are reframed as filter removal rather than signal generation, explaining why DMN suppression yields expanded experience. The receiver cycles at $\approx 40$\,Hz, each cycle a discrete conscious moment; spontaneous thought arises from undirected cycling under internal attractor dynamics. Coupling amplitude decay predicts a pre-registered $\approx 24\%$ reduction in residual gamma power from age 20 to 75 after structural covariates are controlled. Brain death is modelled as receiver failure with the reference pattern continuing; terminal lucidity corresponds to a transient coherence pulse. The framework yields six falsifiable predictions: (1) transient gamma coherence restoration at receiver failure; (2) residual exponential amplitude decay independent of connectome structure; (3) connectome complexity predicting phenomenal richness independently of functional connectivity; (4) a coherence--BOLD dissociation under psilocybin ($\geq 15\%$ coherence increase with $<5\%$ BOLD change in non-DMN regions); (5) a coupling index $A_i(t_n)=\hat{\varphi}^\top\hat{\boldsymbol{\theta}}_i(t_n)$ that outpredicts gamma power for the bispectral index by $\geq 15$ percentage points; and (6) selective recovery of conscious detection under 40\,Hz transcranial entrainment during propofol anaesthesia but not under 10\,Hz control. $\hat{\varphi}$ is estimated once from an independent reference cohort and held fixed; it is not a free parameter. If supported, the model shifts the measurement target from activity magnitude to phase‑structured alignment, with practical consequences for anaesthesia monitoring, covert consciousness detection, and alignment‑sensitive stimulation.
\end{abstract}

\section{Introduction}

Every theory of consciousness begins with the same unresolved problem: why does physical process give rise to subjective experience? The dominant framework---that the brain generates consciousness through neural computation---has produced detailed maps of neural correlates but has not closed the explanatory gap between mechanism and experience \cite{chalmers1995}. We do not dispute the neural correlate data. We dispute the causal direction.

The generator model assumes: neural activity $\rightarrow$ consciousness. We propose the receiver model: $\varphi \rightarrow R_i \rightarrow \varphi_i \rightarrow$ experience. The brain does not produce the signal. It decodes it.

This reframing is motivated by two anomalies that the generator model accounts for only through separate, post-hoc mechanisms. We do not claim these anomalies are inexplicable under the generator view---disinhibition models, network metastability, and terminal neural excitation all offer partial accounts. We claim the receiver model covers both within a single formal structure with no per-anomaly free parameters, and that this unification generates novel, quantitative predictions the generator model does not make.

\textbf{Anomaly 1: Psychedelic expansion under suppression.}
Classical psychedelics (psilocybin, DMT, LSD) dramatically suppress activity in the default mode network \cite{carhart-harris2012}---whose suppression the generator model associates with reduced self-referential processing and ego dissolution. Yet subjective experience does not diminish. It expands, often reported as the most vivid conscious state the subject has encountered. A generator with reduced activity should produce reduced output. It does not. The receiver model predicts this naturally: suppressing the DMN removes a filtering layer, allowing more of $\varphi$ through.

\textbf{Anomaly 2: Terminal lucidity.}
Patients with severe dementia---whose neural tissue is extensively damaged---sometimes recover full, coherent consciousness in the hours before death \cite{nahm2012}. A generator model requires intact hardware for high-quality output. A receiver model allows one final strong coupling between $\varphi$ and the receiver before hardware failure terminates reception entirely.

We do not claim these two anomalies prove the receiver model. We claim they constitute genuine predictive failures of the generator model that motivate a formal alternative. This paper provides that alternative: a phase-alignment coupling framework for consciousness as transmitted, received, and cycled---with equations, predictions, and proposed experiments.

\textbf{Organisation.} Section~2 defines the reference pattern $\hat{\varphi}$ and its properties. Section~3 introduces the receiver $R_i$ and the coupling operator producing individual consciousness $\varphi_i$. Section~4 derives the plant/destroy cycle and its correspondence to neural oscillations. Section~5 models coupling amplitude decay and aging. Section~6 addresses brain death and terminal lucidity. Section~7 proposes six falsifiable experiments. Section~8 discusses scope and open questions.

\section{The Universal Field \texorpdfstring{$\varphi$}{phi}}

\subsection{Definition and Core Properties}

We define $\varphi$ as a scalar field present throughout spacetime, satisfying the following four properties:

\begin{definition}[Universality]
$\varphi$ permeates all matter and space without exclusion. There is no region, medium, or material substrate that attenuates or blocks $\varphi$.
\end{definition}

\begin{definition}[Nonlocality]
$\varphi$ is not spatially bounded. Its state at any point is not a function solely of local conditions. This property is analogous to the nonlocal correlations observed in quantum entanglement \cite{bell1964} but distinct: $\varphi$ is not entangled pairwise between particles but is a globally coherent background field.
\end{definition}

\begin{definition}[Thermal independence]
$\varphi$ does not couple to thermal energy. Its intensity is not a function of temperature. This distinguishes it from thermodynamic quantities and from electromagnetic radiation, both of which are temperature-dependent. The implication is that $\varphi$ does not dissipate with heat loss and is not detectable by calorimetric means.
\end{definition}

\begin{definition}[Receiver-contingent manifestation]
$\varphi$ does not produce experience directly. In the absence of a structured receiver $R_i$, $\varphi$ remains latent---present but unmanifested as subjective experience. Consciousness arises only at the receiver interface, not from the field in isolation.
\end{definition}

\subsection{Formal Representation}

We represent $\varphi$ as a vector of unit norm in gamma-phase space:
\begin{equation}
    \varphi \;\in\; \mathbb{S}^{n_\theta - 1} \;\subset\; \mathbb{R}^{n_\theta},
    \qquad \|\varphi\| = 1
\end{equation}
where $n_\theta$ is the number of EEG source populations. $\varphi$ is the field's \emph{characteristic phase signature}---a reference pattern of gamma-band phase relationships that is globally constant across spacetime (Definition~1).

\textbf{Estimation of $\varphi$.} $\varphi$ is not a free parameter. It is estimated once from a reference cohort and held fixed for all subsequent predictions:
\begin{equation}\label{eq:phi_est}
    \hat{\varphi} \;=\; e_1\!\left(\,\overline{\mathcal{C}}\,\right), \qquad
    \overline{\mathcal{C}} = \frac{1}{N_{\mathrm{ref}}}
    \sum_{k=1}^{N_{\mathrm{ref}}} \mathcal{C}(R_k)
\end{equation}
where $e_1(\cdot)$ denotes the leading eigenvector, $\overline{\mathcal{C}}$ is the mean gamma coherence matrix averaged over a reference cohort of $N_{\mathrm{ref}} \geq 50$ healthy young adults (ages 20--30), and $\mathcal{C}(R_k)$ is each individual's coherence matrix (Section~3.2). The leading eigenvector of the population-averaged coherence is the dominant common gamma-phase pattern shared across healthy young brains---the empirical candidate for the field's phase signature. $\hat{\varphi}$ is unique up to sign (a convention), computed once, and shared across all predictions. All model parameters are empirically estimable from independent datasets and are not tuned to fit outcome variables. Given EEG data and a fixed reference cohort used to estimate $\hat{\varphi}$, all quantities are computable without task-specific fitting.

\textbf{Stability of $\hat{\varphi}$.}
To ensure that $\hat{\varphi}$ is not an artifact of sampling variability or preprocessing choices, we require that it be stable under cohort resampling. Let the reference cohort be partitioned into two disjoint subsets of equal size, producing estimates $\hat{\varphi}_1$ and $\hat{\varphi}_2$. We require:
\begin{equation}
    \left| \hat{\varphi}_1^\top \hat{\varphi}_2 \right| \geq 0.9
\end{equation}
This criterion ensures that the estimated field direction is reproducible across independent samples. Failure to meet this condition constitutes empirical falsification of the field estimate.

We do not write a propagation equation for $\varphi$: the nonlocal character of the field (Definition~2) is inconsistent with a local partial differential equation such as $\Box\varphi = 0$, and the receiver dynamics developed in Sections~3--5 are independent of $\varphi$'s propagation structure. The relevant temporal scale---the plant/destroy cycle frequency---is constrained empirically via correspondence to gamma oscillations in Section~4, not by a dispersion relation.

\subsection{Relation to Known Fields}

$\varphi$ is not proposed as a modification of the Standard Model fields. It is a separate ontological stratum. The analogy is gravitational: gravity was formalised as a field before its source (spacetime curvature) was understood. Newton did not require a mechanism for action at a distance to write $F = Gm_1m_2/r^2$; the predictive structure was established independently. We adopt the same strategy: $\varphi$'s properties are defined operationally through the receiver dynamics they imply, and its ultimate source is treated as an open empirical question.

This stance is not an evasion. It is a deliberate separation of levels: a theory can be well-specified, falsifiable, and useful before its deepest foundations are resolved. Newtonian gravity was all three for 230 years before general relativity supplied the source mechanism. We expect $\varphi$ to follow the same trajectory.

\subsection{What \texorpdfstring{$\varphi$}{phi} Is Not}

Several existing constructs share surface similarities with $\varphi$ but are formally distinct:

\textbf{Integrated Information ($\Phi$).} Tononi's $\Phi$ \cite{tononi2004} is a measure computed over a physical system's causal structure---it quantifies how much information a system generates above its parts. $\varphi$ is not a measure; it is a transmitted field. $\Phi$ is substrate-dependent; $\varphi$ is substrate-independent.

\textbf{Quantum consciousness (Orch-OR).} Penrose and Hameroff \cite{penrose1994} locate consciousness in quantum gravity effects within microtubules. This remains a generator model---the brain produces the event. $\varphi$ does not require quantum gravity and is not localised to any cellular structure.

\textbf{Electromagnetic fields.} Neural EM fields \cite{mcfadden2002} are generated by neural current and are thermally coupled, spatially local, and blocked by standard shielding. All three properties contradict Definitions 2--3 above.

None of these alternatives address both anomalies in Section~1 within a single formal structure.

\subsection{Epistemic Status: What \texorpdfstring{$\varphi$}{phi} Is Not Claiming to Be}

$\varphi$ is \textbf{not} a physical field in the technical sense used in classical or quantum field theory. It is not a solution to a partial differential equation, carries no gauge symmetry, and has no specified propagation dynamics. The term ``field'' is used colloquially to denote a background reference element in $\mathbb{R}^{n_\theta}$ that is present at all receiver locations. The framework is more accurately characterised as an \emph{abstract signal-receiver coupling theory}: its empirical content resides entirely in the measurable quantities $A_i(t_n) = \varphi^\top\hat{\boldsymbol{\theta}}_i(t_n)$ and $\varphi_i(t_n) = A_i(t_n)\cdot\mathbf{W}(R_i)\varphi$, not in any dynamics of $\varphi$ itself. Specifying how $\varphi$ propagates and what conserves it is left as a direction for future work.

The framework does satisfy one conservation-like constraint from the projection geometry. Since $|A_i(t_n)| = |\varphi^\top\hat{\boldsymbol{\theta}}_i| \leq 1$ and $\|\mathbf{W}(R_i)\varphi\| \leq \|\varphi\| = 1$ (projection cannot amplify), the individual consciousness signal is bounded:
\begin{equation}\label{eq:bound}
    0 \;\leq\; \|\varphi_i(t_n)\|^2 \;\leq\; 1 \qquad \forall\, i,\, t_n
\end{equation}
No receiver generates a consciousness signal exceeding the field norm. This is not a dynamical conservation law but a geometric constraint that follows necessarily from the inner product and projection structure of Section~3. It makes a testable prediction: EEG power attributed to the coupling signal $\varphi_i$ should be bounded by the reference power of $\varphi$, with deviations indicating model misspecification.

\section{The Receiver and the Coupling Operator}

\subsection{The Receiver \texorpdfstring{$R_i$}{Ri}}

A receiver $R_i$ is any organised physical system capable of coupling to $\varphi$. We define $R_i$ as a state vector partitioned into three empirically accessible components:
\begin{equation}
    R_i(t) = \bigl(\boldsymbol{\theta}(t),\;\mathbf{w}(t),\;\mathbf{c}(t)\bigr)
    \;\in\; \mathbb{R}^{n_\theta + n_w + n_c}
\end{equation}
where $\boldsymbol{\theta}(t)\in\mathbb{R}^{n_\theta}$ encodes instantaneous oscillation phases of neural populations (extractable from EEG source reconstruction, e.g.\ via Hilbert-transform phase at the gamma band), $\mathbf{w}(t)\in\mathbb{R}^{n_w}$ encodes synaptic weight matrix entries (structural connectivity from DTI fractional anisotropy; functional connectivity from BOLD cross-correlations), and $\mathbf{c}(t)\in\mathbb{R}^{n_c}$ encodes neuromodulator concentrations (dopamine, serotonin, acetylcholine; measurable via PET ligand binding or cerebrospinal fluid assay). These three components capture the receiver's current oscillatory configuration, fixed architecture, and biochemical context---the three parameters that jointly determine which projection of $\varphi$ is accessed at each plant cycle. The specific dimensionalities $n_\theta$, $n_w$, $n_c$ are empirically determined; the framework requires only that $R_i(t)$ be measurable in principle at each $t_n$.

Receivers are not limited to biological brains. The definition is structural: any matter with $C(R) > 0$ constitutes a receiver of non-zero quality, since $Q = 1 - e^{-\beta C} > 0$ for all $C > 0$. Biological nervous systems occupy the high end of this spectrum; their relevance here is empirical, not definitional.

\begin{definition}[Receiver quality]
The quality $Q(R_i) \in [0, 1]$ of receiver $R_i$ is a monotone function of its structural complexity:
\begin{equation}
    Q(R_i) = 1 - e^{-\beta \, C(R_i)}
\end{equation}
where $\beta > 0$ is a scaling constant and $C(R_i)$ is the graph-theoretic entropy of the connectivity subvector $\mathbf{w}$:
\begin{equation}
    C(R_i) \;\equiv\; -\sum_{j} \bar{w}_j \log \bar{w}_j, \qquad
    \bar{w}_j = \frac{|w_j|}{\sum_k |w_k|}
\end{equation}
where $w_j$ are the entries of $\mathbf{w}$ and $\bar{w}_j$ is the normalised edge-weight distribution of the connectome. $C$ is therefore a function of $\mathbf{w}$ alone and is measurable from DTI tractography or BOLD connectivity matrices. As $C \to 0$ (uniform or absent connectivity), $Q \to 0$; as $C \to \infty$, $Q \to 1$.
\end{definition}

Receiver quality governs fidelity, not access: every organised system receives $\varphi$, but high-$Q$ receivers produce richer, more differentiated $\varphi_i$.

\subsection{The Coupling Operator \texorpdfstring{$H$}{H}}

The central operation of the framework is:
\begin{equation}
    \varphi_i(t) = A(\tau)\cdot H\!\bigl(\varphi,\, R_i(t)\bigr)
\end{equation}

$H$ is a coupling operator mapping the universal field $\varphi$ and receiver state $R_i(t)$ to an individual consciousness signal $\varphi_i(t)$. We define $H$ to satisfy four structural constraints:

\textbf{(H1) Receiver-determinism.} For fixed $\varphi$ and fixed age $\tau$, $H$ is a deterministic function of $R_i$. Two receivers in identical states at the same age produce identical $\varphi_i$. (The coupling amplitude $A(\tau)$ scales the output and varies with $\tau$; H1 holds conditionally on $\tau$.)
\begin{equation}
    R_i = R_j,\;\tau_i = \tau_j \implies H(\varphi, R_i, \tau_i) = H(\varphi, R_j, \tau_j)
\end{equation}

\textbf{(H2) Sensitivity.} Small perturbations in $R_i$ can produce distinct $\varphi_i$: $H$ is not globally Lipschitz in $R_i$. There exists a constant $c > 0$ such that for any $R_i$ and any $\delta > 0$ there exist $R_j$ with $\|R_i - R_j\| < \delta$ yet output distance at least $c$:
\begin{equation}
    \exists\, c > 0:\;\forall R_i,\;\forall \delta > 0,\;
    \exists\, R_j:\;
    \|R_i - R_j\| < \delta \;\wedge\;
    \|H(\varphi, R_i) - H(\varphi, R_j)\| \geq c
\end{equation}
Even arbitrarily close receiver states can yield well-separated outputs.

\textbf{(H3) Non-invertibility.} $\varphi_i$ alone does not recover $\varphi$. The full field is not accessible from any single receiver's output; individual experience is a compressed projection, not a lossless copy.

\textbf{(H4) Quality-scaling.} The richness of $\varphi_i$ scales with $Q(R_i)$. Let $\varphi_i^*$ denote the signal produced by a perfect receiver ($Q = 1$). Then:
\begin{equation}
    \mathrm{dim}(\varphi_i) \approx Q(R_i) \cdot \mathrm{dim}(\varphi_i^*)
\end{equation}
where $\mathrm{dim}(\cdot)$ denotes the effective dimensionality of the signal. High-quality receivers access proportionally more dimensions of $\varphi$ than low-quality receivers, whose outputs are coarser projections of the same field.

\textbf{Canonical instantiation.} Because $\varphi \in \mathbb{R}^{n_\theta}$ and $\boldsymbol{\theta}_i(t) \in \mathbb{R}^{n_\theta}$ share the same space, coupling strength is the cosine similarity:
\begin{equation}\label{eq:coupling}
    A_i(t_n) \;=\; \langle \varphi,\, \hat{\boldsymbol{\theta}}_i(t_n) \rangle
    \;=\; \varphi^\top \hat{\boldsymbol{\theta}}_i(t_n)
    \;\in\; [-1, 1]
\end{equation}
where $\hat{\boldsymbol{\theta}}_i(t_n) = \boldsymbol{\theta}_i(t_n) /
\|\boldsymbol{\theta}_i(t_n)\|$ is the normalised Hilbert-transform gamma-phase vector at plant time $t_n$, and $\|\varphi\| = 1$ by definition. The projection matrix is:
\begin{equation}
    \mathbf{W}(R_i) \in \mathbb{R}^{d\times n_\theta}, \qquad
    d = \lceil Q(R_i)\cdot n_\theta \rceil
\end{equation}
whose rows are the top-$d$ eigenvectors of the mean gamma-coherence matrix:
\begin{equation}
    \mathcal{C}(R_i) \;\in\; \mathbb{R}^{n_\theta\times n_\theta}, \qquad
    [\mathcal{C}]_{jk} = \Bigl|\bigl\langle e^{i\theta_j(t)},\, e^{i\theta_k(t)}
    \bigr\rangle_t\Bigr|
\end{equation}
the mean phase coherence between EEG source populations $j$ and $k$, computed from $\boldsymbol{\theta}(t)$. The \textbf{fully computable} expression for individual consciousness at each plant time is:
\begin{equation}\label{eq:main}
\boxed{
    \varphi_i(t_n) \;=\;
    \underbrace{\varphi^\top \hat{\boldsymbol{\theta}}_i(t_n)}_{A_i(t_n)}
    \cdot\;
    \underbrace{\mathbf{W}(R_i(t_n))\,\varphi}_{\text{projection onto coherence eigenmodes}}
}
\end{equation}

\textbf{Model restriction.}
For all predictions and analyses in this paper, we restrict the coupling operator $H$ to the canonical instantiation defined in Eq.~(\ref{eq:main}). Alternative forms of $H$ that satisfy constraints H1--H4 are not considered here. This restriction ensures that the model is uniquely specified and strictly falsifiable. All right-hand quantities are measurable or model parameters: $\boldsymbol{\theta}_i(t_n)$ from EEG Hilbert-transform phases at 40\,Hz plant times; $\mathbf{W}(R_i)$ from EEG-derived gamma coherence eigenvectors; $\varphi \in \mathbb{S}^{n_\theta-1}$ is the fixed reference direction estimated from an independent reference cohort (Section~2.2) and is not adjusted to fit any outcome data. Given data, you can compute $\varphi_i(t_n)$ numerically.

H1 is satisfied by the deterministic map $\boldsymbol{\theta}_i \mapsto \mathcal{C}(R_i) \mapsto \mathbf{W}(R_i)$; H2 by discontinuous eigenvector rotation near eigenvalue degeneracies (generic in empirical coherence matrices); H3 by rank reduction $d < n_\theta$; H4 by the monotone dependence of $d$ on $Q(R_i)$.

\subsection{Individual Consciousness as a Projection}

$\varphi_i$ is the individual's conscious experience. It is neither $\varphi$ itself (which is universal and undifferentiated) nor purely a product of $R_i$ (which is physical structure without phenomenal content). It is the output of their interaction.

This resolves the combination problem that plagues panpsychism: the framework does not require micro-experiences in individual particles to ``combine'' into a unified field. $\varphi$ is already unified; receivers do not combine it, they project from it. The unity of conscious experience is a property of $\varphi$, not an emergent achievement of neural binding.

It also provides a natural account of self-continuity. H2 establishes that $H$ is globally non-Lipschitz: arbitrary receiver states can yield separated outputs. But the actual trajectory $R_i(t)$ is not arbitrary---it is constrained by continuous physical dynamics and stays within a bounded attractor region under normal waking conditions. Along such trajectories, consecutive receiver states $R_i(t_n)$ and $R_i(t_{n+1})$ are close, and $H$ restricted to this region behaves smoothly, so consecutive $\varphi_i$ values are similarly close. The felt sense of an unbroken stream follows from trajectory regularity, not from any global continuity of $H$. Discontinuous changes to $R_i$---anaesthesia, dreamless sleep, severe seizure---push the trajectory outside its normal attractor region, producing large jumps in the output sequence.

\subsection{On Circularity and Direction-Specificity}

A potential objection: since gamma oscillations already correlate with consciousness in the existing literature, using $\boldsymbol{\theta}_i$ (gamma phase) as input to compute $A_i = \hat{\varphi}^\top\hat{\boldsymbol{\theta}}_i$ may appear circular---``consciousness arises from the thing already correlated with consciousness.'' This objection applies to any model that uses gamma as an input variable, but the receiver model is exposed to it in a specific way and must address it directly.

The distinguishing claim is \emph{direction-specificity}. The receiver model does not predict that more gamma power $= $ more consciousness. It predicts that coupling depends on \emph{alignment} between the gamma-phase pattern and the reference direction $\hat{\varphi}$: two receivers with identical gamma power but different phase patterns have different $A_i$. Concretely, tonic-clonic seizures produce high gamma power but grossly disrupted phase coherence; the receiver model predicts $A_i \approx 0$ (random $\hat{\boldsymbol{\theta}}_i$ gives $\hat{\varphi}^\top \hat{\boldsymbol{\theta}}_i \approx 0$ in expectation), consistent with the empirical loss of consciousness. A pure gamma-power model predicts increased consciousness during seizure---which is observationally false.

The psychedelic anomaly (Anomaly~1) is explained within the canonical projection: psilocybin-induced DMN suppression changes $\mathcal{C}(R_i)$, rotating the coherence eigenmodes and increasing $d = \lceil Q(R_i)\cdot n_\theta\rceil$---the number of dimensions of $\hat{\varphi}$ that $\mathbf{W}(R_i)$ projects onto. More dimensions of the reference pattern become accessible, increasing the richness of $\varphi_i$, without any additional free parameters.

\subsubsection*{Edge Case: High Power, Low Alignment}

A core separation from magnitude‑based accounts is the following. If phase structure is random, then
\begin{equation}
    \mathbb{E}[A_i(t)] \approx 0,
\end{equation}
even if gamma power is high. Thus, elevated gamma power is not sufficient for strong conscious output. This yields a direct empirical edge case: states with vigorous neural activity but degraded phase structure (e.g., certain seizures, burst‑suppression patterns) should exhibit weak conscious decoding, a prediction that a pure power‑based view cannot make.

\section{The Plant/Destroy Cycle and Neural Oscillations}

\subsection{Discrete Receiver Dynamics}

The receiver $R_i$ is not a static structure. It is continuously re-instantiated at a characteristic frequency $f$. Each iteration consists of two phases:

\textbf{Plant.} The receiver achieves sufficient structural coherence for $H$ to operate. The computation $\varphi_i = H(\varphi, R_i)$ executes; a moment of conscious experience is produced.

\textbf{Destroy.} The receiver's coherent state dissolves. $H$ is suspended. $\varphi$ continues to transmit but no individual signal is produced until the next plant.

This defines the fundamental quantum of conscious experience: one cycle of duration $T = 1/f$. The subjective sense of continuous experience is not evidence of a continuous process---it is the perceptual integration of rapid discrete cycles, analogous to the apparent continuity of film from discrete frames.

\subsection{Correspondence to Gamma Oscillations}

The plant/destroy cycle frequency $f$ is empirically constrained by the neural correlate literature. Crick and Koch \cite{crick1990} identified 40\,Hz gamma oscillations as the primary candidate for the neural correlate of consciousness---a finding replicated extensively across species, modalities, and paradigms \cite{engel2001}. We interpret gamma oscillations not as generators of consciousness but as the physical signature of the plant/destroy cycle: each gamma oscillation peak marks a plant event; each trough marks a destroy.

\begin{proposition}[Cycle-gamma correspondence]
The plant/destroy cycle frequency satisfies $f \approx 40\,\mathrm{Hz}$, giving a cycle period:
\begin{equation}
    T = \frac{1}{f} \approx 25\,\mathrm{ms}
\end{equation}
Each 25\,ms window constitutes one complete conscious moment: a single evaluation of $H(\varphi, R_i(t))$.
\end{proposition}

This proposition is testable. If the cycle-gamma correspondence holds, then perturbations that disrupt gamma coherence---not merely gamma power---should disrupt conscious experience, while preserved gamma coherence should track preserved consciousness even under other forms of neural disruption. This prediction diverges from the generator model, which requires both coherence and activity to sustain experience.

\subsection{Formal Model}

Let $t_n = nT$ for $n \in \mathbb{Z}_{\geq 0}$ denote the sequence of plant times. The stream of conscious experience is the discrete sequence:
\begin{equation}
    \Phi_i = \left\{\varphi_i(t_n)\right\}_{n=0}^{\infty}, \qquad \varphi_i(t_n) = H\!\left(\varphi,\, R_i(t_n)\right)
\end{equation}

Between plant times---the destroy interval---the receiver is offline. We write:
\begin{equation}
    \varphi_i(t) = \varnothing \qquad \forall\, t \in (t_n,\, t_{n+1})
\end{equation}

The perceived continuity of experience arises from perceptual integration across consecutive cycles, not from $H$ operating continuously. The visual flicker-fusion threshold ($\approx$60\,Hz) is the frequency \textit{above} which discrete stimuli are perceptually fused; below it, flickering can in principle be detected. The plant cycle at 40\,Hz falls below this threshold, but the threshold applies to external stimuli presented to the perceptual system---the plant cycle is the perceptual system itself, not an input to it. Temporal binding operates over the cycle outputs rather than the substrate, producing the felt sense of continuity. Film at 24\,fps is the appropriate analogy: continuity arises from integration over discrete frames, not from the frames themselves being continuous.

\subsection{Sleep, Anaesthesia, and Cycle Disruption}

The cycle model predicts that altered states of consciousness correspond to disruptions of $f$, not simply reductions in neural activity. We distinguish three modes:

\textbf{Normal waking.} $f \approx 40$\,Hz, cycles regular, $R_i$ coherent at each plant time. Full $\varphi_i$ stream produced.

\textbf{Dreamless sleep / deep anaesthesia.} Gamma coherence collapses; plant events fail to complete. The receiver attempts to cycle but cannot achieve the coherence threshold for $H$ to execute. Output: $\varphi_i = \varnothing$ across extended intervals. This is loss of consciousness, not loss of $\varphi$.

\textbf{REM sleep / dreaming.} Gamma coherence is partially restored under reduced sensory input. $H$ executes but with $R_i$ driven by internal state rather than external stimulation, producing $\varphi_i$ with high vividness but poor external correspondence.

The anaesthesia case is particularly important. Anaesthetic agents are known to disrupt thalamo-cortical gamma synchrony specifically \cite{alkire2008}. The receiver model predicts this is the mechanism of unconsciousness: thalamo-cortical disruption prevents the plant phase from completing. The generator model cannot easily distinguish why \textit{this specific frequency band} is necessary if the brain is simply computing experience from activity.

\subsection{Spontaneous Thought as Undirected Cycling}

When external input is absent, the receiver continues to cycle at $f$. The coupling operator still executes: $\varphi_i(t_n) = H(\varphi, R_i(t_n))$. But $R_i$ is no longer shaped by sensory input; it is shaped by internal attractor dynamics. The output $\varphi_i$ is determined by whichever attractor state $R_i$ inhabits at each $t_n$.

This provides a mechanistic account of spontaneous thought and mind-wandering---the target phenomenon of many consciousness studies. Spontaneous thought is not noise. It is the $\varphi_i$ stream produced when the receiver cycles without external constraint, allowing internal attractor states to dominate $R_i$ and thus $H(\varphi, R_i)$. Different attractors (memory, emotion, imagination) produce qualitatively different projections of $\varphi$, which manifest as different categories of spontaneous mental content.

\section{Coupling Amplitude Decay and Aging}

\subsection{Coupling Amplitude as an Inner Product}

The coupling amplitude at each plant cycle is the inner product defined in Eq.~(\ref{eq:coupling}):
\begin{equation}
    A_i(t_n) = \varphi^\top\hat{\boldsymbol{\theta}}_i(t_n) \;\in\; [-1,1]
\end{equation}
This quantity fluctuates cycle-to-cycle. Its expected value across cycles at organismal age $\tau$ defines the \emph{mean coupling amplitude}:
\begin{equation}
    \bar{A}(\tau) = \mathbb{E}_{t_n}\!\bigl[\varphi^\top\hat{\boldsymbol{\theta}}_i(t_n)\bigr]
\end{equation}
$\bar{A}(\tau)$ is directly measurable: record high-density EEG, extract Hilbert-transform gamma-phase vectors $\boldsymbol{\theta}_i(t_n)$ at each 40\,Hz cycle, project onto $\varphi$, and average over the recording epoch. As gamma coherence degrades with age, $\hat{\boldsymbol{\theta}}_i(t_n)$ rotates away from the fixed reference direction $\varphi$, reducing $\bar{A}(\tau)$ independently of the structural term $Q$.

Whereas $Q$ measures \emph{fidelity} (dimensionality of the projection), $A_i$ measures \emph{coupling strength} (instantaneous phase-alignment with $\varphi$). A high-$Q$ receiver with misaligned gamma produces a rich but quiet $\varphi_i$; a low-$Q$ receiver strongly aligned produces a strong but coarse signal.

\subsection{Exponential Amplitude Decay}

We model the mean coupling amplitude as an exponentially decaying function of organismal age $\tau$:
\begin{equation}
    \bar{A}(\tau) = A_0\, e^{-\lambda \tau}
\end{equation}
where $A_0 = \bar{A}(0) > 0$ is the mean coupling at birth and $\lambda > 0$ is the decay constant. This functional form is a \emph{hypothesis}, not an axiom: it predicts how $\mathbb{E}_{t_n}[\varphi^\top\hat{\boldsymbol{\theta}}_i]$ behaves across the lifespan. It is motivated by the empirical finding that gamma coherence degrades with age through myelin loss, synaptic pruning, and reduced neurovascular efficiency; the exponential form arises under the assumption that this misalignment grows at a constant rate $\lambda$. Prediction~2 (Section~7) tests this directly.

The amplitude enters the coupling operator as a scaling factor:
\begin{equation}
    \varphi_i(\tau) = A(\tau)\cdot H(\varphi,\, R_i)
    = A_0\, e^{-\lambda \tau} \cdot H(\varphi,\, R_i)
\end{equation}

As $\tau$ increases, $\varphi_i$ is attenuated even when receiver structure $R_i$ is held constant.

\subsection{Correspondence to Neural Oscillation Slowing}

A well-documented empirical regularity is that peak gamma frequency decreases with age in healthy humans \cite{polich2007}. The generator model treats this as a consequence of neural degradation reducing computational throughput. The receiver model offers a distinct account: as $A(\tau)$ decays, the driving force sustaining each plant/destroy cycle weakens.

Specifically, we propose that the plant phase at each discrete time $t_n$ requires the per-cycle coupling to exceed a threshold:
\begin{equation}
    \text{plant succeeds at }t_n \iff A_i(t_n)\cdot Q(R_i(t_n)) \geq A_{\min}, \quad A_{\min}\in(0,1)
\end{equation}
where $A_i(t_n)=\hat{\varphi}^\top\hat{\boldsymbol{\theta}}_i(t_n)\in[-1,1]$ is the per-cycle cosine similarity (not the population mean $\bar{A}(\tau)$). Using the population mean in the threshold condition would conflate a smooth age-level statistic with a cycle-level gate; the condition must act on the instantaneous realisation.

The effective cycle rate is the fraction of cycles that clear the threshold:
\begin{equation}
    f_{\mathrm{eff}}(\tau) \;=\;
    f_{\max}\cdot\Pr_{t_n}\!\bigl[A_i(t_n)\cdot Q(R_i)\geq A_{\min}\bigr]
    \;\approx\; f_{\max}\cdot\frac{\bar{A}(\tau)}{A_0}
    \;=\; f_{\max}\cdot e^{-\lambda\tau}
\end{equation}
where the approximation holds when the success probability is monotone in the mean coupling $\bar{A}(\tau)=A_0 e^{-\lambda\tau}$. As $\bar{A}(\tau)$ decays with age, fewer cycles clear the threshold and $f_{\mathrm{eff}}$ falls. This predicts that peak gamma frequency follows the same exponential envelope as $\bar{A}(\tau)$---a falsifiable claim (Section~7).

\textbf{Gamma power spectrum prediction.} Because signal amplitude scales as $A(\tau)$, EEG gamma band power (which scales as amplitude squared) satisfies:
\begin{equation}
    S_\gamma(\tau) = A_0^2\, e^{-2\lambda\tau}\cdot Q\bigl(R_i(\tau)\bigr)\cdot S_{\gamma,0}
\end{equation}
After partialing out the structural term $Q(R_i)$ via DTI fractional anisotropy and cortical thickness covariates, the residual gamma power satisfies:
\begin{equation}\label{eq:gamma_decay}
    \tilde{S}_\gamma(\tau) = S_{\gamma,0}\, e^{-2\lambda\tau}
\end{equation}
The decay exponent $2\lambda$ is twice the coupling amplitude decay rate and is structurally independent of $Q$. Calibrating from published data on peak gamma frequency decline ($f_{\mathrm{peak}}\approx 40$\,Hz at age 20, $\approx 35$\,Hz at age 75; see Section~7), $\hat{\lambda}\approx 0.0025\,\mathrm{yr}^{-1}$, giving $2\hat{\lambda}\approx 0.005\,\mathrm{yr}^{-1}$. Equation~(\ref{eq:gamma_decay}) then predicts a $1 - e^{-0.005\times 55}\approx 24\%$ reduction in gamma power from age 20 to 75 that survives structural covariate control. The generator model predicts that once $Q$-related structural covariates are controlled, no systematic residual trend should remain.

\subsection{Aging Phenomenology}

The amplitude decay model makes contact with reported phenomenology of aging that the generator model has not systematically addressed:

\textbf{Reduced vividness.} Older adults consistently report that experiences feel less vivid and intense than in youth \cite{mather2012}. Under the generator model this is attributed to reduced neural activity. Under the receiver model it is a direct prediction: $A(\tau)$ decays, so $\varphi_i$ is attenuated regardless of $Q(R_i)$. A well-preserved elderly brain still receives a quieter signal.

\textbf{Slower processing.} Cognitive processing speed declines monotonically with age. The receiver model attributes this to lower $f_{\mathrm{eff}}$: fewer completed cycles per second means fewer evaluated moments of $\varphi_i$, producing slower integration of information.

\textbf{Sleep degradation.} REM duration and gamma coherence during sleep both decline with age \cite{mander2017}. Under the receiver model: the already-reduced $A(\tau)$ makes sustaining the coherence threshold during sleep---when external input is absent to support $R_i$---progressively harder.

Note that $Q(R_i)$ also declines with age due to structural neural degradation. The two effects are additive: amplitude decay attenuates the signal while receiver degradation reduces fidelity. Disentangling the two contributions is an empirical task; Section~7 proposes experimental handles.

\subsection{Childhood as Peak Amplitude}

The model implies that $\tau = 0$ corresponds to maximum amplitude $A_0$. Early childhood, before receiver $Q$ has been fully optimised by learning and pruning, involves a high-amplitude but lower-fidelity signal: a strong $\varphi_i$ passed through a still-developing filter stack. This accounts for well-documented features of childhood cognition---heightened emotional intensity, vivid imagination, lower resistance to unusual percepts---without positing any special neural mechanism. The signal is simply louder.

\section{Death, Dispersal, and Reconcentration}

\subsection{Brain Death as Receiver Failure}

Within the receiver model, brain death is the permanent termination of $R_i$---the receiver fails and $H$ can no longer execute. This is categorically different from the termination of $\varphi$, which continues to transmit irrespective of any individual receiver's state. The field does not switch off; the decoder does.

Formally, let $\tau_{\dagger}$ denote the time of receiver failure. For $\tau > \tau_{\dagger}$:
\begin{equation}
    R_i(\tau) = \varnothing \implies \varphi_i(\tau) = H(\varphi, \varnothing) = \varnothing
\end{equation}

The individual conscious signal ceases. $\varphi$ is unaffected.

\subsection{Terminal Lucidity}

Before examining post-death dynamics, we address the empirical anomaly of terminal lucidity (Section~1, Anomaly~2). Under the receiver model, terminal lucidity is a consequence of a transient coherence event in a degrading receiver.

As neurodegeneration advances, $Q(R_i)$ declines monotonically. Most cycles fail to complete because $A(\tau) \cdot Q(R_i) < A_{\min}$: this is the mechanism of dementia-related unconsciousness. However, physical systems near catastrophic failure sometimes undergo a final structural reorganisation---a transient return to coherence before terminal collapse, analogous to a critically loaded system passing through a metastable state.

We model this as a terminal coherence pulse of finite duration $\sigma > 0$ (corresponding to the observed minutes-to-hours window of lucidity events):
\begin{equation}
    Q(R_i, \tau) = Q_{\mathrm{base}}(\tau)
    + \Delta Q\cdot\exp\!\left(-\frac{(\tau - (\tau_{\dagger} - \sigma))^2}{2\sigma^2}\right)
\end{equation}
where $\Delta Q > 0$ is the peak magnitude and $\sigma$ is the characteristic duration. The transient increase $\Delta Q$ is not a free postulate: excitotoxic ionic cascade at receiver failure triggers coordinated mass depolarisation across neural populations, transiently increasing phase coherence $\mathcal{C}(R_i)$ and therefore $Q$ before terminal collapse \cite{borjigin2023}. $\Delta Q$ is thus a predicted consequence of the same synchronisation event independently documented in cardiac arrest patients; its magnitude and $\sigma$ are constrained by the observed minutes-to-hours window of terminal lucidity events \cite{nahm2012}. If $A_i(t_n)\cdot(Q_{\mathrm{base}}+\Delta Q)\geq A_{\min}$ during the pulse, plant phases complete: $H$ executes and a final vivid $\varphi_i$ is produced.

This predicts that terminal lucidity events should be accompanied by a transient restoration of gamma coherence immediately before death---a testable claim that can be evaluated with continuous EEG monitoring in palliative patients.

\section{Falsifiable Predictions and Proposed Experiments}

The receiver model makes predictions that diverge from the generator model in directions accessible to current or near-term technology. We organise them by the section of theory they test.

\subsection{Prediction 1: Terminal Gamma Coherence Restoration}

\textbf{Derivation.} Section~6 predicts a terminal coherence spike in $Q(R_i)$ immediately before receiver failure, producing a brief window of complete cycle execution and a final vivid $\varphi_i$.

\textbf{Experiment.} Continuous high-density EEG monitoring of palliative patients with severe dementia or other terminal neurodegeneration. Measure 40\,Hz band coherence (not power) as a function of time-to-death, computed retrospectively.

\textbf{Receiver model prediction.} A statistically significant transient increase in gamma coherence in the final minutes to hours before clinical brain death, co-occurring with reported terminal lucidity events.

\textbf{Generator model prediction.} Terminal EEG changes are expected as a non-specific death artifact: excitotoxic depolarisation and ionic cascade at cell death can produce transient high-frequency bursts \cite{borjigin2023}. However, these surges are artifacts of shutdown mechanics and should not be correlated with subjective lucidity reports---the brain is incapable of generating coherent experience when tissue is irreversibly damaged.

\textbf{Distinguishing feature.} Both models can produce a terminal gamma signal; they differ on its relationship to lucidity. The receiver model predicts gamma coherence restoration is \textit{specifically} and \textit{temporally} locked to reported lucidity events. The generator model predicts any terminal coherence change is a non-specific artifact with no correlation to subjective reports. The test is therefore: compute the temporal offset between gamma coherence peaks and documented lucidity windows. A significant correlation falsifies the generator account; absence of correlation falsifies the receiver account.

\subsection{Prediction 2: Exponential Gamma Frequency Decay with Age}

\textbf{Derivation.} Section~5 predicts $f_{\mathrm{eff}}(\tau) = f_{\max}\cdot A(\tau)/A_0 \propto e^{-\lambda\tau}$ in the above-threshold regime. Because $A(\tau)$ is a coupling property of the receiver--field interface rather than of neural structure, the exponential trend should persist after structural degradation metrics are fully controlled.

\textbf{Experiment.} Longitudinal resting-state EEG study with repeated measures across decades (or a large cross-sectional design, $n \geq 30$ per decade from age 10 to 80, with cohort-effects controls). Extract peak gamma frequency per subject. Regress out cortical thickness, white matter integrity (DTI fractional anisotropy), and synaptic density (PET-based) to partial out $Q(R_i)$. Fit both a linear and an exponential decay model to the residuals.

\textbf{Receiver model prediction.} After controlling for all available structural metrics, peak gamma frequency residuals follow an exponential decay with age with rate $\hat{\lambda} = 0.0025\,\mathrm{yr}^{-1}$ (calibrated from published frequency data; Section~5). The gamma band power residuals satisfy $\tilde{S}_\gamma(\tau) = S_{\gamma,0}e^{-2\hat{\lambda}\tau}$, predicting a $1 - e^{-2\times 0.0025\times 55} = 24\%$ reduction from age 20 to 75. The pre-registered falsification criterion is: a one-tailed likelihood-ratio test comparing the exponential to a linear age trend in the covariate-partialled residuals, $\alpha = 0.05$, $n \geq 30$ per decade. Under $\hat{\lambda} = 0.0025\,\mathrm{yr}^{-1}$, the predicted effect size is Cohen's $d \approx 0.6$ for the exponential vs.\ linear contrast---detectable with the specified $n$. If the exponential fit is not significantly better than linear after structural covariate control, the coupling amplitude model is falsified.

\textbf{Generator model prediction.} Frequency decline is fully explained by structural degradation metrics; once $Q$-related covariates are controlled, no systematic age trend remains in the residuals. A linear model fits as well or better than exponential.

\textbf{Distinguishing feature.} Residual exponential trend ($d \geq 0.6$) surviving structural covariate control, at the pre-registered $\hat{\lambda}$. A trend that exists but at a different rate than $\hat{\lambda}$ would also falsify the model as specified.

\subsection{Prediction 3: Connectome Complexity and Phenomenal Richness}

\textbf{Derivation.} Section~3 defines $Q(R_i) = 1 - e^{-\beta C(R_i)}$, where $C(R_i)$ is structural complexity. Higher $Q$ predicts richer $\varphi_i$---more dimensions of $\varphi$ projected into experience.

\textbf{Experiment.} Acquire whole-brain structural connectivity via DTI and resting-state functional connectivity via fMRI. Compute both graph-theoretic entropy $C(R_i)$ of the structural connectome and a standard functional integration measure (e.g., global efficiency of the BOLD correlation matrix). Administer validated phenomenological richness measures (Vividness of Visual Imagery Questionnaire, experience sampling of spontaneous thought complexity, visual masking thresholds). Test whether structural $C(R_i)$ predicts phenomenal richness after controlling for functional connectivity, BOLD signal amplitude, and IQ.

\textbf{Receiver model prediction.} Structural $C(R_i)$ explains variance in phenomenal richness \textit{beyond} functional connectivity---because $Q$ sets the dimensionality of $\varphi_i$ through structural architecture, independent of moment-to-moment functional dynamics. A receiver with high structural complexity but currently suppressed functional connectivity (e.g., during low-arousal rest) should still exhibit richer phenomenology than a low-complexity receiver.

\textbf{Generator model prediction.} Phenomenal richness tracks functional connectivity (moment-to-moment integration); structural complexity explains no additional variance once functional measures are controlled, because what matters for generation is active computation, not static wiring.

\textbf{Distinguishing feature.} Structural complexity predicting richness \textit{independently of} functional connectivity. Both models predict a structural--richness correlation; only the receiver model predicts it survives functional covariate control.

\subsection{Prediction 4: Psychedelic Coherence--BOLD Decoupling}

\textbf{Derivation.} Section~3 explains psilocybin effects through the canonical projection: DMN suppression alters $\mathcal{C}(R_i)$, increasing the accessible dimensionality $d$. This predicts a specific dissociation absent from generator accounts.

\textbf{Prior evidence.} Carhart-Harris et al.\ \cite{carhart-harris2012} already documented DMN suppression under psilocybin and increased cross-network integration, providing preliminary support. The experiment proposed here targets the dose--response relationship, which constitutes an independent test not present in the 2012 data.

\textbf{Experiment.} Simultaneous fMRI and EEG across a graded psilocybin dose series (e.g., 0, 5, 10, 20\,mg oral) in a within-subjects design. At each dose, quantify DMN suppression via BOLD and non-DMN 40\,Hz coherence via EEG.

\textbf{Receiver model prediction.} At the highest dose (20\,mg), mean fronto-parietal 40\,Hz coherence should increase by $\geq 15\%$ from placebo baseline during the peak-effect window (60--120\,min post-dose), while BOLD in non-DMN regions changes by $< 5\%$ from baseline. Specifically: if $\bar{\gamma}$ denotes mean non-DMN 40\,Hz coherence across parieto-frontal electrodes,
\begin{equation}
    \frac{\bar{\gamma}_{20\,\mathrm{mg}} - \bar{\gamma}_{0\,\mathrm{mg}}}
         {\bar{\gamma}_{0\,\mathrm{mg}}} \;\geq\; 0.15
\end{equation}
The derivation: each additional milligram removes further filter attenuation; the $\geq 15\%$ threshold is calibrated from the Carhart-Harris 2012 cross-network integration increase ($\sim$20\% by their metric), projected conservatively to coherence. The monotone dose--response should hold across all four doses (Jonckheere--Terpstra trend test, $\alpha = 0.05$). If the 20\,mg condition shows $< 15\%$ coherence increase, or if the BOLD constraint is violated (non-DMN BOLD increases $\geq 5\%$), the filter-removal model is falsified.

\textbf{Generator model prediction.} Coherence increase should scale with non-DMN BOLD ($r \geq 0.5$ expected); if coherence rises while BOLD is flat, the generator model lacks a mechanism.

\textbf{Distinguishing feature.} $\geq 15\%$ coherence increase at 20\,mg \textit{combined with} $< 5\%$ BOLD change in non-DMN regions. Either criterion alone is insufficient; both must hold simultaneously.

\subsection{Prediction 5: Coupling Index Outpredicts Gamma Power for BIS}

\textbf{Derivation.} The receiver model's core claim is direction-specificity: consciousness depends on \emph{alignment} $A_i(t_n) = \hat{\varphi}^\top \hat{\boldsymbol{\theta}}_i(t_n)$, not on gamma power alone. This yields a concrete, conservative falsification test against existing monitoring standards.

\textbf{Experiment.} During gradual propofol induction (effect-site concentration stepped from 0 to $4\,\mu\mathrm{g/mL}$ in $0.5\,\mu\mathrm{g/mL}$ increments, $n \geq 20$ subjects), record 64-channel EEG continuously. At each 40\,Hz plant time $t_n$: compute (a) $A_i(t_n) = \hat{\varphi}^\top\hat{\boldsymbol{\theta}}_i(t_n)$ using the pre-estimated reference $\hat{\varphi}$; (b) gamma band power $P_\gamma(t_n)$; (c) gamma coherence $\bar{C}_\gamma(t_n)$. Regress BIS index (recorded simultaneously) on each predictor separately, then compare $r^2$ values.

\textbf{Receiver model prediction.}
\begin{equation}
    r^2(A_i \to \mathrm{BIS}) \;\geq\; 0.60
    \quad\text{and}\quad
    r^2(A_i \to \mathrm{BIS}) - r^2(P_\gamma \to \mathrm{BIS}) \;\geq\; 0.15
\end{equation}
The coupling index $A_i$ should explain at least 15 percentage points more BIS variance than gamma power, because it captures the direction of phase alignment that power discards. This gap should be largest in conditions where power is preserved but alignment is disrupted (e.g., at moderate propofol concentrations where gamma persists but consciousness degrades).

\textbf{Generator model prediction.} Gamma power (or coherence) predicts BIS as well as $A_i$; the partial $r^2$ of $A_i$ after controlling for $P_\gamma$ is $\leq 0.05$.

\textbf{Falsification criterion.} If $r^2(A_i) - r^2(P_\gamma) < 0.15$ across all subjects ($\alpha = 0.05$, one-tailed paired test), the direction-specificity claim is falsified. This experiment does not require novel equipment---only a reference EEG dataset and the pre-estimated $\hat{\varphi}$---and is replicable on any existing depth-of-anaesthesia EEG dataset.

\subsection{Prediction 6: 40\,Hz Transcranial Entrainment During Propofol Anaesthesia}

\textbf{Derivation.} If the receiver plant/destroy cycle is genuinely tied to gamma‑band phase alignment, then externally entraining the cortex at 40\,Hz should selectively facilitate conscious processing, whereas entrainment at other frequencies should not.

\textbf{Experiment.} Healthy volunteers ($n \geq 20$) undergo gradual propofol sedation while 64‑channel EEG is recorded. At a steady‑state effect‑site concentration where behavioural responsiveness is absent (e.g., $2.5\,\mu\mathrm{g/mL}$), apply three tACS conditions in randomised blocks: (i) 40\,Hz fronto‑parietal entrainment, (ii) 10\,Hz control entrainment, and (iii) sham. During each block, administer a forced‑choice auditory‑visual stimulus and measure the stimulus‑locked gamma coherence and the behavioural detection rate.

\textbf{Receiver model prediction.}
\begin{itemize}
    \item 40\,Hz entrainment will increase stimulus‑locked 40\,Hz coherence by $\geq 20\%$ relative to sham and raise detection rate from $\approx 0\%$ to $\geq 30\%$.
    \item 10\,Hz entrainment will produce no significant change in detection rate ($< 5\%$ improvement) despite equal current density.
    \item The coupling index $A_i(t_n)$ computed from the entrained EEG will show a significant increase over sham, and this increase will mediate the detection improvement (mediation model: $c'$ path non‑significant).
\end{itemize}

\textbf{Generator model prediction.} Entrainment at any frequency may alter cortical excitability, but no frequency‑specific recovery of consciousness is expected beyond what overall arousal modulation would produce. Any behavioural improvement will be independent of gamma phase alignment with the fixed $\hat{\varphi}$.

\textbf{Falsification criterion.} If 40\,Hz entrainment fails to produce a behavioural recovery significantly greater than 10\,Hz entrainment (mixed‑effects logistic regression, interaction $p < 0.05$), the gamma‑cycle hypothesis is falsified. This experiment does not require novel recording equipment; it adds only a tACS stimulator to a standard anaesthesia EEG setup.

\section{Technological and Measurement Relevance}

If the receiver framework is supported, it suggests that consciousness may be more accurately indexed by structured phase alignment than by spectral power alone. This would motivate new measurement approaches in at least four areas:
\begin{enumerate}
    \item \textbf{Clinical monitoring of awareness during anaesthesia.} Current depth‑of‑anaesthesia monitors (e.g., BIS) rely heavily on power and entropy measures. An alignment‑based index, computed as $A_i = \hat{\varphi}^\top \hat{\boldsymbol{\theta}}_i$, could provide a more direct readout of conscious access, particularly in states where gamma power is preserved but coherence is lost.
    \item \textbf{Detection of covert consciousness in behaviourally unresponsive patients.} Patients with disorders of consciousness may retain residual phase structure even when motor output is absent. Alignment metrics could offer a non‑invasive window into the presence of a functioning receiver.
    \item \textbf{Improved decoding of conscious state transitions.} Sleep‑wake boundaries, emergence from anaesthesia, and seizure‑related loss of consciousness are characterised by abrupt changes in coherence organisation. Tracking alignment with a fixed reference $\hat{\varphi}$ may reveal these transitions earlier and more reliably than power‑based markers.
    \item \textbf{Development of alignment‑sensitive stimulation paradigms.} Closed‑loop stimulation that adjusts its phase to maximise $\hat{\varphi}^\top \hat{\boldsymbol{\theta}}_i$ could be used therapeutically to maintain or restore conscious processing in fragile receivers.
\end{enumerate}
The value of the framework is therefore not just conceptual; it changes what counts as the relevant variable of measurement.

\section{Scope, Limitations, and Open Questions}

\subsection{What This Theory Does and Does Not Claim}

The signal theory addresses the \textit{causal direction} of the relationship between brain and consciousness. It does not deny that neural correlates exist; it reinterprets them as receiver dynamics rather than generator outputs. Every fMRI and EEG finding in the existing literature is compatible with the receiver model---the data constrain the receiver, not the field.

The theory does not solve the hard problem in the sense of deriving subjective experience from physical substrate. What it does is relocate the hard problem: rather than asking how neural computation becomes experience, we ask how $\varphi$ carries experiential content.

\textbf{What this model uniquely predicts.} Three predictions distinguish the receiver model from all existing theories in ways that are testable with current technology: (1) Gamma band spectral power decays as $e^{-2\lambda\tau}$ \emph{after} structural covariates are controlled (Prediction~2)---no existing theory predicts a residual exponential trend structurally independent of the connectome; (2) Coupling index $A_i = \hat{\varphi}^\top\hat{\boldsymbol{\theta}}_i$ predicts BIS with $r^2 \geq 0.60$, exceeding gamma power by $\geq 15$ percentage points (Prediction~5)---no existing theory predicts this direction-specific advantage; (3) The psychedelic coherence--BOLD decoupling at $\geq 15\%$ with $< 5\%$ BOLD change (Prediction~4)---disinhibition and metastability accounts do not predict the BOLD constraint. These three predictions are pre-registered, numerically specific, and obtainable with existing EEG/fMRI infrastructure.

\subsection{Agency and Phenomenological Structure}

One testable instance of the coupling framework is felt agency---the experience of deliberation, volition, and choice. These are well-localised to prefrontal cortical circuits, which under the receiver model constitute a filter layer $F_{\mathrm{PFC}}$ shaping the projection of $\varphi$ into experience that includes agency-phenomenology. Two predictions follow:

\textbf{(i) Prefrontal damage preferentially reduces experienced volition.} Damage to $F_{\mathrm{PFC}}$ should attenuate agency-content in $\varphi_i$ preferentially over other phenomenal dimensions---consistent with documented avolition as a leading symptom in prefrontal lesion patients, though prefrontal damage also affects attention and working memory, reflecting the multiple filter functions of this region.

\textbf{(ii) The Libet readiness potential is not paradoxical.} Neural preparation precedes reported conscious intention by $\approx$550\,ms \cite{libet1983}. Generator models require additional assumptions---such as a retrospective conscious veto---to reconcile this with felt agency. The receiver model predicts it directly: the 550\,ms window spans approximately 22 plant cycles ($550\,\mathrm{ms}/25\,\mathrm{ms}$). During this window $R_i$ undergoes preparatory structural reconfiguration---the motor program is being instantiated in $R_i$---and the plant threshold $A(\tau)\cdot Q(R_i)\geq A_{\min}$ is only met once that reconfiguration reaches sufficient coherence. The intention emerges in $\varphi_i$ only when $R_i$ first achieves the coherence threshold for the motor-state configuration. The experience of choosing is the decoded output of a multi-cycle physical process, appearing in $\varphi_i$ only at its conclusion.

The framework suggests, as a working hypothesis, that qualitative phenomenological structure more broadly---felt continuity, unity, selfhood---may similarly reflect specific filter configurations in $H$. This is not derived from H1--H4 and constitutes an open research direction rather than a current prediction. The source of $\varphi$ itself remains an open empirical question, bounded by what receiver measurements can constrain.

\subsection{Non-Human Receivers}

The receiver definition (Section~3) is structural, not biological. Any organised system with $C(R) > 0$ constitutes a receiver of quality $Q(R) > 0$. This implies a gradient of conscious experience across all organised matter, from negligible (simple systems) to maximal (complex biological brains). It does not imply that rocks are conscious in any phenomenologically meaningful sense; it implies they are receivers of negligibly small quality.

\subsection{Relationship to Existing Frameworks}

The receiver model is not irreconcilable with existing theories:

\textbf{Global Workspace Theory \cite{baars1988}.} The global workspace can be reinterpreted as the mechanism by which $R_i$ achieves the structural coherence required for plant-phase completion. Broadcast across the workspace corresponds to the receiver reaching the coherence configuration at which $H$ can execute, not to clearing the amplitude threshold $A_{\min}$ per se.

\textbf{Predictive Processing \cite{clark2016}.} The hierarchical prediction machinery is a model of how $R_i$ is structured and updated. The receiver model is agnostic about the internal architecture of $R_i$; predictive processing specifies one plausible implementation.

\textbf{Integrated Information Theory.} IIT's $\Phi$ is a candidate for $C(R_i)$: a high-$\Phi$ system has high structural complexity and would thus have high $Q$. The receiver model and IIT are compatible at the level of receiver quality, but diverge fundamentally on causal direction: IIT holds that $\Phi$ generates experience; the receiver model holds that $\Phi$ determines fidelity of reception.

\textbf{Recurrent Processing Theory \cite{crick1990,engel2001}.} Lamme's recurrent processing theory holds that reentrant feedback between cortical areas is the neural correlate of consciousness---feedforward-only processing is unconscious; recurrence is what generates experience. The receiver model is compatible with this finding while reversing its interpretation: recurrent feedback is the mechanism by which $R_i$ sustains the phase coherence required for $H$ to execute. Disrupting recurrence reduces $Q(R_i)$ and therefore $A_i(t_n)$, predicting loss of conscious access---the same empirical result, derived from receiver dynamics rather than generation. RPT describes receiver architecture; it does not identify the source of experience.

\textbf{Higher-Order Theories.} Higher-order theories (Rosenthal) hold that a mental state is conscious only when accompanied by a higher-order representation of that state. Within the receiver framework, higher-order representations correspond to the top eigenvectors selected by $\mathbf{W}(R_i)$: the projection matrix selects the dimensions of $R_i$ that carry integrated, self-referential structure, and only those dimensions participate in the coupling. A receiver lacking higher-order representations has a degenerate $\mathbf{W}$ and low $Q$, yielding low $A_i$. Higher-order structure is thus a contributor to receiver quality, not the generator of experience itself.

\subsection{Limitations}

The framework has three acknowledged limitations:

\textbf{(L1) Underdetermination of $H$.} The coupling operator is constrained by four structural requirements but not uniquely specified. Multiple functional forms satisfy H1--H4. Identifying $H$ empirically is a long-term research programme, not a near-term prediction.

\textbf{(L2) Reference cohort sensitivity.} $\hat{\varphi}$ is estimated as the leading eigenvector of $\overline{\mathcal{C}}$ pooled from a reference cohort of $N \geq 50$ healthy adults (ages 20--30). The estimate is stable when the cohort is well-matched, but will shift if the cohort includes atypical EEG profiles (e.g.\ high rates of subclinical epilepsy or medication use). Robustness checks across independent cohorts and pre-registration of cohort-selection criteria are required before the predictions in Section~7 can be evaluated without circularity.

\textbf{(L3) Free scaling parameters.} Two threshold parameters remain uncalibrated: $\beta > 0$ in $Q(R_i)=1-e^{-\beta C(R_i)}$ and $A_{\min}\in(0,1)$ in the per-cycle triggering condition. Both are in principle identifiable---$\beta$ from the empirical relationship between DTI-derived connectivity entropy and phenomenal richness ratings; $A_{\min}$ from the coupling amplitude at which conscious detection falls to chance in the psychophysics literature. Until calibrated, predictions that depend on the threshold (e.g.\ $f_{\mathrm{eff}}$ and the terminal lucidity pulse condition) are qualitative rather than quantitative.

\textbf{(L4) No quantum account.} The framework is written at the level of classical field theory. Whether $\varphi$ requires a quantum description---and what that description would be---is not addressed. This is a gap, not a refutation: classical electromagnetism was complete and useful before quantum electrodynamics.

\subsection{Conclusion}

We have presented a field-theoretic framework in which consciousness is transmitted through organised matter rather than generated by it. The universal field $\varphi$ interacts with receivers via a coupling operator $H(\varphi, R_i(t))$, where $R_i(t) = (\boldsymbol{\theta}(t), \mathbf{w}(t), \mathbf{c}(t))$ is a measurable state vector encoding oscillation phases, synaptic weights, and neuromodulator concentrations. The receiver cycles at $f \approx 40$\,Hz; coupling amplitude $A(\tau) = A_0 e^{-\lambda\tau}$ decays exponentially with age, predicting gamma band spectral power to fall as $e^{-2\lambda\tau}$ after structural covariates are controlled---a pre-registered, quantitative prediction with $\hat{\lambda}\approx 0.0025\,\mathrm{yr}^{-1}$ and an estimated 24\% power reduction from age 20 to 75. At death the receiver fails while $\varphi$ continues; terminal lucidity is modelled as a transient coherence pulse in the degrading receiver immediately before failure.

The framework is falsifiable at six experimental junctures identified in Section~7, culminating in the propofol/tACS entrainment experiment (Prediction~6) which provides a clear disconfirmation pathway for the gamma-as-plant-phase hypothesis. The framework reframes two anomalies that the generator model cannot cleanly explain, subsumes existing theories as special cases of receiver dynamics, and relocates the hard problem without claiming to dissolve it. The source of $\varphi$ remains an open question---as it should, at this stage of the theory.

\bibliographystyle{plain}
\bibliography{signal_theory}

\end{document}